\section{Requerimiento de Control Integral}

Para determinar si es necesario aplicar una acción de control integral, se debe evaluar el error en estado estacionario frente a la topología y precisión del sistema físico implementado.

Desde el punto de vista teórico, el sistema propuesto \textbf{no requiere} control integral. Esto se justifica al analizar el Tipo de Sistema de la planta a lazo abierto. Como se determinó previamente en la Ecuación (1), la función de transferencia $G_{(s)} = \frac{G_0}{s^2}$ posee dos polos en el origen, lo que la clasifica como un sistema de \textbf{Tipo 2} (contiene dos integradores puros en cascada). La teoría de control establece que cualquier sistema con al menos un integrador en la trayectoria directa garantiza por naturaleza un error en estado estacionario nulo ante una entrada escalón. Matemáticamente, agregar una acción integral adicional resulta redundante.

En la implementación práctica, se registró que la tensión de estabilización alcanzó $\sim 0.98\text{V}$ frente a la referencia escalón de $1.0\text{V}$ (ver Tabla \ref{tab:params_tiempo}). Sin embargo, este error residual del $\sim 2\%$ no representa una incapacidad del lazo de control para seguir la referencia, sino que es un \textbf{error de escalamiento en la etapa de realimentación}. 

Físicamente, la ganancia $k_1$ se implementó mediante resistores. Debido a las tolerancias comerciales de los componentes reales, la ganancia de medición física difiere ligeramente del $k_1 = 1,00$ ideal. La planta cumplió perfectamente su función integrando la señal hasta que el error detectado en el nodo sumador fue cero, pero debido a esta desviación física en $k_1$, el sistema se enclavó en un valor de $y = \frac{1\text{V}}{k_{1_{real}}} \approx 0,98\text{V}$.

\textbf{Conclusión:} En el contexto de este diseño analógico, la aplicación de control integral es innecesaria y desaconsejable por los siguientes motivos:
\begin{itemize}
	\item \textbf{Redundancia y Estabilidad:} Añadir un tercer integrador a la trayectoria directa degradaría severamente la estabilidad relativa del sistema (reduciendo el margen de fase), lo que dificultaría seriamente el cumplimiento de la restricción de sobrepico máxima solicitada.
	\item \textbf{Ineficacia ante el offset físico:} Un control integral adicional integraría el error basándose en la misma medición realimentada. Al estar la ganancia $k_1$ desviada por tolerancias físicas de los resistores, el integrador extra forzaría al sistema a estabilizarse exactamente en el mismo valor de $0.98\text{V}$.
	\item \textbf{Complejidad de Hardware:} Requeriría incorporar una etapa adicional con otro amplificador operacional y red pasiva, aumentando el costo y la complejidad del circuito sin aportar ningún beneficio real al seguimiento de la referencia.
\end{itemize}